\documentclass[a4paper,11pt]{article}
\usepackage[it]{regodoc}
\RdColor{azure}{2F5F8A}
\RdColor{slate}{4A5560}
\RdColor{sage}{4F6F4B}
\RdColor{bronze}{7D5F2A}
\RdColor{plum}{674766}
\RdColor{brick}{8E4436}
\RdColor{olive}{5F6328}
\RdColor{gray}{5E5E5E}
\RdColor{ok}{2F6B3A}
\RdColor{warning}{9A6200}
\RdColor{critical}{A12B2B}
\RdColor{muted}{6E6E6E}
\RdNames{Parte}{Figura}{Indice}
\RdTitle{Rinnovo della rete della Cantina Colle Rovere}
\RdSubject{Analisi dello stato attuale e proposta}
\RdReference{CR-2026-07}
\RdVersionLabel{Bozza 0.3}
\RdDate{5 ottobre 2026}
\begin{document}
\RdTitlePage
\RdTableOfContents
\begin{RdLegend}{Legenda delle etichette}
\RdLegendItem{slate}{osservato}
\RdLegendItem{warning}{da verificare}
\end{RdLegend}

\RdOpening{rd-summary}

\RdSection{1}{rd-in-brief}
{In breve}
{In breve}

Il rapporto riguarda la rete e gli archivi della cantina; i lavori proposti partono
il 16 novembre 2026 e richiedono al massimo 14~giorni. Domande e osservazioni vanno
all’Ufficio tecnico della cantina.
\par
\begin{RdList}{numbered}
\RdItem Rifare la copertura senza fili della barricaia.
\RdItem Affidare la manutenzione a un installatore, per 320–410~euro l’anno.
\RdItem Provare il recupero dei file a ogni fine vendemmia.
\end{RdList}

Le analisi complete sono in \RdRef{rd-network}{Parti~I–II}.
\par

\RdPart{I}{azure}{rd-network}
{Rete e collegamenti}
\RdSummary{Rete degli uffici, collegamento a Internet e copertura senza fili della cantina.}

\RdSection{1}{rd-lan}
{Rete locale}
{Rete locale}

\RdLabel{slate}{osservato} Le \RdMono{14} postazioni di ufficio e laboratorio passano
da un apparato centrale di \RdI{vecchia generazione}; il dominio interno è \RdDomain{colle-\RdAB{}rovere.\RdAB{}example}.
\par

Le segnalazioni vanno a \RdMail{6D61696C746F3A697440636F6C6C652D726F766572652E6578616D706C65}{it@\RdAB{}colle-\RdAB{}rovere.\RdAB{}example};
la guida dei punti di accesso è in \RdLink{68747470733A2F2F646F63732E6578616D706C652E636F6D2F61702F696E7374616C6C617A696F6E65}{Installazione
dei punti di accesso}\RdCite{rd-ap-guide}{1}.
\par

\RdSection{2}{rd-wifi}
{Copertura senza fili}
{Copertura senza fili}

La barricaia è coperta al 72,5\%: oltre la seconda fila di botti il segnale è \RdTone{warning}{debole},
una “zona d’ombra” per i lettori delle etichette.
\par
\RdFigure{images/coverage.png}{0.6}
{Copertura misurata prima dei lavori.}

\RdSection{1}{rd-sources-network}
{Fonti}
{Fonti}
\begin{RdSources}{Fonti}
\RdSource{1}{rd-ap-guide}{68747470733A2F2F646F63732E6578616D706C652E636F6D2F61702F696E7374616C6C617A696F6E65}
{https:\RdAB{}/\RdAB{}/\RdAB{}docs.\RdAB{}example.\RdAB{}com/\RdAB{}ap/\RdAB{}installazione}
{Guida all’installazione dei punti di accesso}
{Example Networks}
{consultata il 10 settembre 2026}
\RdSource{2}{}{68747470733A2F2F646F63732E6578616D706C652E636F6D2F6E61732F636F706965}
{https:\RdAB{}/\RdAB{}/\RdAB{}docs.\RdAB{}example.\RdAB{}com/\RdAB{}nas/\RdAB{}copie}
{Copie pianificate sul NAS}
{Example Storage}
{}
\end{RdSources}

\RdPart{II}{sage}{rd-backup}
{Copie di sicurezza}
\RdSummary{Dove stanno gli archivi della cantina e come renderli più sicuri.}

\RdSection{1}{rd-current-backup}
{Archivi e copie}
{Archivi e copie}

\RdLabel{slate}{osservato} Il disco di rete in ufficio è quasi pieno; il nuovo
disco di copia avrà \RdMono{4~TB}, e dopo un guasto il recupero completo richiederà
circa 6~ore.
\par

Per le copie fuori dalla cantina vedi \RdRef{rd-offsite}{“Copia esterna”}.
\par
\begin{RdBox}{sage}{}

Funziona già bene:
\par
\begin{RdList}{bullet}
\RdItem le versioni precedenti dei fogli di calcolo;
\RdItem gli accessi separati per ufficio e laboratorio.
\end{RdList}
\end{RdBox}
\begin{RdTable}{normal}{3}{1}{4.8cm}
{>{\RdCellStart}p{\RdStar}>{\RdCellStart}p{2.2cm}>{\RdCellStart}p{2.6cm}}
\RdCell{head}{1}{\RdStar}{Archivio}
&\RdCell{head}{1}{2.2cm}{Dove}
&\RdCell{head}{1}{2.6cm}{Copia}
\RdHeadEnd
\RdCell{row}{1}{\RdStar}{Registro di cantina e analisi di laboratorio}
&\RdCell{row}{1}{2.2cm}{disco di rete}
&\RdCell{row}{1}{2.6cm}{\RdTone{warning}{settimanale}}
\RdRowEnd
\RdCell{row}{1}{\RdStar}{Etichette e grafica:\par
\begin{RdList}{bullet}
\RdItem file di stampa;
\RdItem fotografie delle annate.
\end{RdList}}
&\RdCell{row}{1}{2.2cm}{portatile grafico}
&\RdCell{row}{1}{2.6cm}{\RdTone{critical}{nessuna} \RdLabel{warning}{da verificare}}
\RdRowEnd
\RdSeparator
\RdCell{total}{2}{\dimexpr\RdStar+2.2cm+2\tabcolsep\relax}{Capacità del nuovo disco}
&\RdCell{total}{1}{2.6cm}{\RdMono{4} TB}
\RdRowEnd
\end{RdTable}

\RdSection{2}{rd-offsite}
{Copia esterna}
{Copia esterna}

Una copia lontana dalla cantina protegge anche da incendio e furto. La prova a
secco del trasferimento:
\par
\begin{RdListing}
\RdListingLine{rclone~sync~-{}\RdAB{}-\RdAB{}dry-\RdAB{}run~/\RdAB{}volume1/\RdAB{}archivi/~esterno:\RdAB{}colle-\RdAB{}rovere/}
\RdListingLine{echo~\textquotedbl{}prova~conclusa\textquotedbl{}}
\end{RdListing}
\begin{RdSmall}

Con \RdCode{-{}\RdAB{}-\RdAB{}dry-\RdAB{}run} il comando elenca i file senza trasferirli.
\par
\end{RdSmall}

\RdCheckContent
\end{document}
